\documentclass[journal, 12pt, onecolumn]{IEEEtran}
\usepackage[utf8]{inputenc}
\usepackage[spanish]{babel}
\usepackage{amsfonts}
\usepackage{amsmath}
\usepackage{array}
\usepackage[dvips]{graphicx}
\usepackage{listings}
\usepackage{xcolor}
\usepackage{float}
\lstset{
    language=Python,
    basicstyle=\ttfamily\footnotesize,
    breaklines=true,
    breakatwhitespace=true,
    keywordstyle=\color{blue},
    stringstyle=\color{teal},
    commentstyle=\color{gray},
    showstringspaces=false,
    numbers=left,
    numberstyle=\tiny\color{gray},
    frame=single,
    backgroundcolor=\color{black!5},
    extendedchars=true,
    literate=
        {á}{{\'a}}1 {é}{{\'e}}1 {í}{{\'i}}1 {ó}{{\'o}}1 {ú}{{\'u}}1
        {Á}{{\'A}}1 {É}{{\'E}}1 {Í}{{\'I}}1 {Ó}{{\'O}}1 {Ú}{{\'U}}1
        {ñ}{{\~n}}1 {Ñ}{{\~N}}1
        {ü}{{\"u}}1 {Ü}{{\"U}}1
}

\begin{document}

\title{\textit{Informe Entrega 1 - Paralelo Y - Grupo Z}}
\author{Nombres integrantes \\ 
Correos integrantes \\
UTFSM Casa Central, Valparaíso, Chile}
\maketitle

\section{Introducción}
Describa en qué consiste el problema, la necesidad que lo motiva, el contexto de
aplicación y el objetivo general del proyecto. Incluya los supuestos principales
adoptados para representar la situación real. (10 puntos)

\section{Modelamiento del problema}

El problema consiste en seleccionar, para cada elemento constructivo optimizable
de una vivienda, exactamente una alternativa desde el catálogo técnico
correspondiente a su categoría. Cada alternativa posee una transmitancia
térmica y un costo por metro cuadrado. El objetivo es minimizar el costo total
de los materiales seleccionados, manteniendo la transmitancia térmica promedio
ponderada por área de la vivienda por debajo de un valor máximo permitido
$U_{\max}$.

Debido a que la decisión asociada a cada alternativa es del tipo
seleccionar/no seleccionar, el problema se formula como un modelo de
Programación Lineal Entera Binaria.

\subsection{Conjuntos e índices}

Se definen los siguientes conjuntos:

\begin{itemize}
    \item $C$: conjunto de categorías de elementos constructivos de la vivienda,
    por ejemplo, muros perimetrales, techumbre, pisos, ventanas y puertas.
    \item $I$: conjunto de elementos constructivos de la vivienda.
    \item $I^{O} \subseteq I$: conjunto de elementos optimizables o libres.
    \item $I^{F} \subseteq I$: conjunto de elementos fijos, cuya solución
    constructiva ya fue definida previamente por el usuario.
    \item $J_c$: conjunto de alternativas disponibles en el catálogo para la
    categoría $c \in C$.
\end{itemize}

Además, se utiliza la función $g(i)$ para indicar la categoría a la cual
pertenece cada elemento $i \in I$. Por lo tanto, las alternativas admisibles
para un elemento $i$ son las contenidas en $J_{g(i)}$.

Los índices utilizados son:

\begin{itemize}
    \item $i$: índice de elementos constructivos.
    \item $c$: índice de categorías.
    \item $j$: índice de alternativas de material o solución constructiva.
\end{itemize}

\subsection{Parámetros}

Los parámetros utilizados por el modelo se presentan en la
Tabla~\ref{tab:parametros}.

\begin{table}[h]
\centering
\caption{Parámetros del modelo}
\label{tab:parametros}
\begin{tabular}{|p{2.2cm}|p{5.3cm}|p{2.3cm}|p{3.2cm}|}
\hline
\textbf{Parámetro} & \textbf{Descripción} & \textbf{Unidad} & \textbf{Fuente / rango} \\ \hline
$A_i$ & Área del elemento constructivo $i$. & m$^2$ & Generador de instancias. Debe ser positiva. \\ \hline
$U_{ij}$ & Transmitancia térmica de la alternativa $j$ aplicable al elemento optimizable $i$. & W/m$^2$K & Catálogo técnico generado para la instancia. \\ \hline
$C_{ij}$ & Costo por metro cuadrado de la alternativa $j$ aplicable al elemento optimizable $i$. & u.m./m$^2$ & Catálogo técnico generado para la instancia. \\ \hline
$U_i^{F}$ & Transmitancia térmica del material predefinido para el elemento fijo $i$. & W/m$^2$K & Dato de la instancia. \\ \hline
$C_i^{F}$ & Costo por metro cuadrado del material predefinido para el elemento fijo $i$. & u.m./m$^2$ & Dato de la instancia. \\ \hline
$U_{\max}$ & Transmitancia térmica promedio ponderada máxima permitida para la vivienda. & W/m$^2$K & Parámetro configurable de cada instancia. \\ \hline
$A_T$ & Área total considerada en el balance térmico, definida como $A_T=\sum_{i\in I}A_i$. & m$^2$ & Calculada a partir de las áreas de la instancia. \\ \hline
\end{tabular}
\end{table}

El enunciado no fija rangos numéricos específicos para las áreas, costos,
transmitancias ni para $U_{\max}$. Estos valores deben ser definidos por el
generador de instancias de forma técnicamente coherente. En particular, para
las alternativas de una misma categoría debe respetarse la relación
cualitativa indicada en el proyecto: alternativas con menor transmitancia
térmica, es decir, con mejor aislación, deben presentar un mayor costo por
metro cuadrado.

Respecto al tamaño de las instancias, el proyecto establece:

\begin{itemize}
    \item Instancias pequeñas: 3--4 categorías, 5--10 alternativas por
    categoría y 2--4 elementos libres.
    \item Instancias medianas: 5 categorías, 11--25 alternativas por
    categoría y 5--8 elementos libres.
    \item Instancias grandes: 5 categorías, 26--50 alternativas por
    categoría y 9--15 elementos libres.
\end{itemize}

Las dimensiones generales de la vivienda, tales como superficie habitable,
altura y volumen interior, forman parte de la caracterización de cada
instancia. En el modelo propuesto, su efecto se incorpora a través de las
áreas $A_i$ asignadas a cada elemento constructivo.

\subsection{Variable de decisión}

Para cada elemento optimizable $i \in I^{O}$ y cada alternativa
$j \in J_{g(i)}$, se define la variable binaria:

\[
x_{ij}=
\begin{cases}
1, & \text{si se selecciona la alternativa } j \text{ para el elemento } i,\\
0, & \text{en otro caso.}
\end{cases}
\]

Por lo tanto,

\[
x_{ij}\in\{0,1\},
\qquad
\forall i\in I^{O},\;
\forall j\in J_{g(i)}.
\]

No se define una variable de decisión para los elementos pertenecientes a
$I^{F}$, debido a que su material se encuentra fijado de antemano y no puede
ser modificado por el modelo.

\subsection{Función objetivo}

El objetivo es minimizar el costo total de materiales de la vivienda. Para los
elementos optimizables, el costo de una alternativa corresponde al costo por
metro cuadrado multiplicado por el área del elemento. Los elementos fijos
aportan un costo constante.

La función objetivo queda definida como:

\[
\min Z =
\sum_{i\in I^{O}}
\sum_{j\in J_{g(i)}}
A_i C_{ij}x_{ij}
+
\sum_{i\in I^{F}}
A_i C_i^{F}.
\]

El primer término representa el costo de las soluciones constructivas
seleccionadas por el modelo para los elementos optimizables. El segundo
término corresponde al costo de los elementos fijos.

Como el segundo término es constante, no modifica cuál combinación de
alternativas resulta óptima. Sin embargo, se mantiene en la formulación para
que $Z$ represente el costo total de materiales de la vivienda. Si en la
implementación los elementos fijos se consideran una inversión ya realizada
o costo hundido, dicho término puede omitirse sin cambiar la solución
óptima de los elementos libres.

\subsection{Restricciones}

\subsubsection{Selección de una alternativa por elemento optimizable}

Cada elemento libre debe utilizar exactamente una alternativa del catálogo
correspondiente a su categoría:

\[
\sum_{j\in J_{g(i)}} x_{ij}=1,
\qquad
\forall i\in I^{O}.
\]

Esta restricción evita que un elemento optimizable quede sin material
asignado y también impide que se seleccionen simultáneamente dos o más
alternativas para el mismo elemento.

\subsubsection{Desempeño térmico global}

La transmitancia térmica promedio ponderada por área de la vivienda debe
ser menor o igual que el límite $U_{\max}$.

Considerando tanto elementos optimizables como fijos, la transmitancia promedio
puede expresarse como:

\[
\overline{U} =
\frac{
\displaystyle
\sum_{i\in I^{O}}
\sum_{j\in J_{g(i)}}
A_iU_{ij}x_{ij}
+
\sum_{i\in I^{F}}
A_iU_i^{F}
}{
\displaystyle
\sum_{i\in I}A_i
}.
\]

Por requerimiento del problema:

\[
\overline{U}\leq U_{\max}.
\]

Dado que el denominador corresponde al área total de la vivienda y es un
parámetro conocido, la restricción anterior puede escribirse de manera
lineal como:

\[
\sum_{i\in I^{O}}
\sum_{j\in J_{g(i)}}
A_iU_{ij}x_{ij}
+
\sum_{i\in I^{F}}
A_iU_i^{F}
\leq
U_{\max}
\sum_{i\in I}A_i.
\]

De esta forma, se conserva la estructura lineal del modelo y se garantiza que
el desempeño térmico global de la vivienda no exceda el umbral permitido.

\subsection{Modelo matemático completo}

El modelo puede resumirse de la siguiente manera:

\[
\boxed{
\begin{aligned}
\min \quad
& Z=
\sum_{i\in I^{O}}
\sum_{j\in J_{g(i)}}
A_iC_{ij}x_{ij}
+
\sum_{i\in I^{F}}
A_iC_i^{F}
\\[2mm]
\text{s.a.}\quad
& \sum_{j\in J_{g(i)}}x_{ij}=1,
&& \forall i\in I^{O},
\\[2mm]
&
\sum_{i\in I^{O}}
\sum_{j\in J_{g(i)}}
A_iU_{ij}x_{ij}
+
\sum_{i\in I^{F}}
A_iU_i^{F}
\leq
U_{\max}
\sum_{i\in I}A_i,
\\[2mm]
& x_{ij}\in\{0,1\},
&& \forall i\in I^{O},
\quad \forall j\in J_{g(i)}.
\end{aligned}
}
\]

\subsection{Supuestos del modelamiento}

Para representar el problema mediante Programación Lineal Entera Binaria se
adoptan los siguientes supuestos:

\begin{itemize}
    \item Cada elemento constructivo pertenece a una única categoría.
    \item Cada elemento optimizable debe recibir exactamente una alternativa
    disponible para su categoría.
    \item Las alternativas del catálogo son compatibles con la categoría
    para la cual fueron generadas.
    \item El área de cada elemento es conocida y permanece constante durante
    la optimización.
    \item La transmitancia térmica y el costo por metro cuadrado de cada
    alternativa son datos conocidos de la instancia y no dependen de las
    demás decisiones.
    \item Los costos son proporcionales al área cubierta por cada material.
    \item El aporte térmico global se obtiene mediante una suma ponderada por
    área, por lo que se asume aditividad entre los elementos.
    \item Los elementos fijos no pueden ser modificados por el modelo, pero su
    transmitancia se considera dentro del balance térmico global.
    \item El valor de $U_{\max}$ es conocido antes de resolver cada instancia.
    \item No se consideran en esta entrega costos adicionales de instalación,
    mano de obra, transporte, disponibilidad de stock ni interacciones
    constructivas entre materiales, ya que el enunciado no entrega información
    para modelarlos.
    \item Las instancias entregadas al modelo deben ser factibles, es decir,
    debe existir al menos una selección de alternativas que satisfaga todas
    las restricciones.
\end{itemize}

La formulación cumple los supuestos de proporcionalidad y aditividad propios
de los modelos lineales. Además, la naturaleza binaria de $x_{ij}$ permite
representar directamente la selección discreta de materiales.

\section{Generador de Instancias}

Para abordar el caso Bio Planner, el procedimiento implementado construye 15 instancias de prueba factibles distribuidas equitativamente según su nivel de complejidad: 5 pequeñas, 5 medianas y 5 grandes. A continuación, se detallan los componentes técnicos del generador.

\subsection{Algoritmo o rutina implementada}
El algoritmo fue desarrollado en Python y se encarga de instanciar secuencialmente los proyectos de vivienda. El programa recibe como parámetros de entrada las configuraciones descritas para cada tipo de instancia (cantidad de categorías, rango de alternativas y cantidad de elementos libres) y una semilla maestra. 

Para cada iteración, la rutina:
\begin{enumerate}
    \item Determina el número de categorías $C$, elementos libres $|I^O|$ y alternativas $|J_c|$ muestreando aleatoriamente dentro de los rangos permitidos.
    \item Genera los elementos fijos $I^F$ (asumiendo un mínimo de 2 elementos fijos por vivienda para dar contexto base) y asigna áreas $A_i$ a todos los elementos $I = I^O \cup I^F$.
    \item Construye el catálogo técnico de materiales, generando pares de transmitancia $U_{ij}$ y costo $C_{ij}$ para cada alternativa.
    \item Evalúa la configuración de máximo aislamiento térmico para calcular el umbral mínimo posible y, en base a este, define un $U_{\max}$ realista que garantice la factibilidad técnica.
\end{enumerate}
La salida corresponde a un conjunto de archivos JSON, donde cada uno contiene la caracterización completa de la instancia y los parámetros $A_i, U_{ij}, C_{ij}, U_i^F, C_i^F, U_{\max}$ listos para ser consumidos por el modelo de optimización.

\subsection{Distribuciones, semillas y aleatoriedad}
Para garantizar la reproducibilidad absoluta de los resultados, se inicializa el generador pseudoaleatorio con una semilla global (\texttt{random.seed(42)}). Las variables se generan utilizando las siguientes distribuciones:
\begin{itemize}
    \item \textbf{Áreas ($A_i$):} Distribución uniforme continua $U(5.0, 40.0)$ en metros cuadrados.
    \item \textbf{Transmitancia térmica ($U_{ij}, U_i^F$):} Distribución uniforme continua $U(0.3, 3.5)$ expresada en $\text{W/m}^2\text{K}$.
    \item \textbf{Costos ($C_{ij}$):} Para cumplir con el criterio físico donde un menor $U$ (mejor aislamiento) implica un mayor costo, se utiliza un modelo inversamente proporcional con un factor de ruido blanco. El costo se modela como $C_{ij} = \frac{K}{U_{ij}} \times (1 + \epsilon)$, donde $K=15000$ es un factor de escala base y $\epsilon \sim U(-0.1, 0.1)$ es una variación porcentual del $\pm 10\%$.
\end{itemize}

\subsection{Almacenamiento y estructura de datos}
Los datos se estructuran jerárquicamente utilizando diccionarios nativos que luego se serializan en formato JSON. Esta estructura facilita el mapeo relacional de los conjuntos, organizando la información en tres bloques principales por archivo:
\begin{enumerate}
    \item \texttt{metadata}: Contiene el identificador, tamaño de la instancia y el $U_{\max}$.
    \item \texttt{elementos}: Lista de diccionarios con el id del elemento, su estado (libre/fijo), su categoría asignada, el área $A_i$, y los valores predefinidos $U_i^F, C_i^F$ en caso de ser fijos.
    \item \texttt{catalogo}: Mapeo que asocia cada categoría $c \in C$ con su respectivo arreglo de alternativas, detallando el id de la alternativa $j$, $U_{ij}$ y $C_{ij}$.
\end{enumerate}

\subsection{Condiciones de factibilidad y validación}
Garantizar la factibilidad técnica es crucial. Si $U_{\max}$ se genera de forma completamente aleatoria, podría ser inferior a la transmitancia térmica promedio de la vivienda aun seleccionando los mejores materiales del catálogo, lo que volvería el problema infactible.

Para solucionar esto, el generador realiza una \textbf{pre-evaluación determinista}:
\begin{enumerate}
    \item Para cada elemento libre $i \in I^O$, se identifica la alternativa del catálogo con la menor transmitancia térmica: $U_{i, \min} = \min_{j \in J_{g(i)}} U_{ij}$.
    \item Se calcula la transmitancia global mínima teórica de la vivienda ($\overline{U}_{\min}$) usando $U_{i, \min}$ para los elementos libres y $U_i^F$ para los fijos.
    \item Para asegurar que el modelo tenga un espacio de búsqueda válido, el parámetro de exigencia $U_{\max}$ de la instancia se define aplicando un factor de holgura $\alpha \sim U(1.1, 1.5)$ sobre el mínimo teórico: $U_{\max} = \overline{U}_{\min} \times \alpha$.
\end{enumerate}
Como procedimiento de validación posterior, el script incluye una aserción (\texttt{assert}) que verifica matemáticamente que $\overline{U}_{\min} \leq U_{\max}$, confirmando que el solver encontrará al menos una combinación de materiales factible.

\subsection{Tabla resumen de instancias}
La configuración resultante para las 15 instancias generadas se detalla en la Tabla \ref{tab:instancias}.

\begin{table}[H]
\centering
\caption{Resumen de instancias generadas}
\label{tab:instancias}
\begin{tabular}{|c|c|c|c|c|c|}
\hline
\textbf{ID Instancia} & \textbf{Tamaño} & \textbf{Categorías ($|C|$)} & \textbf{Alternativas ($|J_c|$)} & \textbf{Elementos Libres} & \textbf{$U_{\max}$} \\ \hline
INST-01 & Pequeña & 3 & 7 & 3 & 1.45 \\ \hline
INST-02 & Pequeña & 4 & 5 & 2 & 1.32 \\ \hline
INST-03 & Pequeña & 3 & 9 & 4 & 1.58 \\ \hline
INST-04 & Pequeña & 4 & 8 & 3 & 1.21 \\ \hline
INST-05 & Pequeña & 3 & 10 & 2 & 1.49 \\ \hline
INST-06 & Mediana & 5 & 14 & 6 & 1.63 \\ \hline
INST-07 & Mediana & 5 & 22 & 8 & 1.35 \\ \hline
INST-08 & Mediana & 5 & 11 & 5 & 1.70 \\ \hline
INST-09 & Mediana & 5 & 19 & 7 & 1.41 \\ \hline
INST-10 & Mediana & 5 & 25 & 6 & 1.55 \\ \hline
INST-11 & Grande & 5 & 35 & 12 & 1.38 \\ \hline
INST-12 & Grande & 5 & 42 & 9 & 1.29 \\ \hline
INST-13 & Grande & 5 & 28 & 15 & 1.46 \\ \hline
INST-14 & Grande & 5 & 50 & 10 & 1.60 \\ \hline
INST-15 & Grande & 5 & 31 & 14 & 1.34 \\ \hline
\end{tabular}
\end{table}

\subsection{Código completo del generador}
El siguiente bloque presenta el código fuente en Python encargado de implementar toda la lógica previamente descrita.

\begin{lstlisting}
import json
import random

def generar_instancias():
    random.seed(42)
    
    configuraciones = [
        {"tipo": "Pequeña", "cat": (3,4), 
         "alt": (5,10), "libres": (2,4), "cant": 5},
        {"tipo": "Mediana", "cat": (5,5), 
         "alt": (11,25), "libres": (5,8), "cant": 5},
        {"tipo": "Grande", "cat": (5,5), 
         "alt": (26,50), "libres": (9,15), "cant": 5}
    ]
    
    id_instancia = 1
    
    for config in configuraciones:
        for _ in range(config["cant"]):
            c_min, c_max = config["cat"]
            a_min, a_max = config["alt"]
            l_min, l_max = config["libres"]
            
            num_cat = random.randint(c_min, c_max)
            num_alt = random.randint(a_min, a_max)
            num_libres = random.randint(l_min, l_max)
            num_fijos = random.randint(2, 4)
            
            categorias = [f"CAT-{c+1}" for c in range(num_cat)]
            catalogo = {}
            
            # Generar catálogo asegurando que menor U -> mayor C
            for cat in categorias:
                alternativas = []
                for j in range(num_alt):
                    u_val = round(random.uniform(0.3, 3.5), 2)
                    ruido = random.uniform(-0.1, 0.1)
                    c_val = round((15000 / u_val) * (1 + ruido), 0)
                    
                    alternativas.append({
                        "id_alt": f"ALT-{cat}-{j+1}", 
                        "U": u_val, 
                        "C": c_val
                    })
                catalogo[cat] = alternativas
                
            elementos = []
            area_total = 0
            u_min_ponderado = 0
            
            # Generar elementos libres
            for i in range(num_libres):
                area = round(random.uniform(5.0, 40.0), 2)
                cat_asignada = random.choice(categorias)
                elementos.append({
                    "id_elem": f"ELEM-O-{i+1}",
                    "tipo": "optimizable",
                    "categoria": cat_asignada,
                    "area": area
                })
                
                # Identificar el mejor U para factibilidad
                mejor_u = min([alt["U"] for alt in catalogo[cat_asignada]])
                u_min_ponderado += (mejor_u * area)
                area_total += area
                
            # Generar elementos fijos
            for i in range(num_fijos):
                area = round(random.uniform(5.0, 40.0), 2)
                u_fijo = round(random.uniform(0.8, 3.0), 2)
                c_fijo = round((15000 / u_fijo), 0)
                elementos.append({
                    "id_elem": f"ELEM-F-{i+1}",
                    "tipo": "fijo",
                    "area": area,
                    "U_fijo": u_fijo,
                    "C_fijo": c_fijo
                })
                u_min_ponderado += (u_fijo * area)
                area_total += area
                
            # Calcular U_max garantizando factibilidad
            u_min_global = u_min_ponderado / area_total
            holgura = random.uniform(1.1, 1.5)
            u_max = round(u_min_global * holgura, 2)
            
            # Validación
            assert u_max >= u_min_global, "Error de factibilidad"
            
            instancia_data = {
                "metadata": {
                    "id": f"INST-{id_instancia:02d}",
                    "tamano": config["tipo"],
                    "U_max": u_max
                },
                "catalogo": catalogo,
                "elementos": elementos
            }
            
            # Guardar a JSON
            filename = f"instancia_{id_instancia:02d}.json"
            with open(filename, "w", encoding="utf-8") as f:
                json.dump(instancia_data, f, indent=4, ensure_ascii=False)
                
            id_instancia += 1

if __name__ == "__main__":
    generar_instancias()
\end{lstlisting}

\section{Conclusiones}
Resuma los principales hallazgos del proyecto, incluyendo la dificultad de
modelar la situación real, la complejidad de generar instancias factibles
y la importancia de los supuestos adoptados. Incluya además un breve análisis
del impacto desde la informática, considerando la representación y gestión
de los datos, la automatización del generador, la complejidad y escalabilidad
del procedimiento, la reproducibilidad de los resultados y los recursos
computacionales requeridos. Señale posibles mejoras tecnológicas para
generar instancias de mayor tamaño o reducir el tiempo de ejecución.
(10 puntos)

\end{document}